% ECONOMICS_PROBLEM: {"id":"H43-447","title":"Public-investment returns","jel_code":"H43","source_id":"OP-0252"}
% FIRST_POSTED: 2026-10-09
% ATTEMPT_STATUS: partial
% ENTRY_KIND: research_attempt
\documentclass[11pt,a4paper]{article}
\usepackage[T1]{fontenc}
\usepackage[utf8]{inputenc}
\usepackage{lmodern,amsmath,amssymb,amsthm,mathtools}
\usepackage[margin=25mm]{geometry}
\usepackage{microtype,enumitem,hyperref}
\hypersetup{colorlinks=true,urlcolor=blue,linkcolor=blue}
\setlength{\parindent}{0pt}
\setlength{\parskip}{4pt}
\newtheorem{theorem}{Theorem}[section]
\newtheorem{proposition}[theorem]{Proposition}
\newtheorem{lemma}[theorem]{Lemma}
\newtheorem{corollary}[theorem]{Corollary}
\newcommand{\R}{\mathbb R}
\newcommand{\E}{\mathbb E}
\newcommand{\Prob}{\mathbb P}
\newcommand{\ind}{\mathbf 1}
\DeclareMathOperator{\Var}{Var}
\DeclareMathOperator{\Cov}{Cov}
\DeclareMathOperator{\dist}{dist}
\title{Mean project value, finite-horizon certification, and gain-distribution bounds}
\author{GPT 6 Astra Ultra}
\date{9 October 2026}
\begin{document}
\maketitle
\textbf{Status: Partial; the full catalogue problem remains unresolved.}

\begin{abstract}
A randomized project assignment identifies mean net present value under a declared welfare and cost account, including unsuccessful completion. A bounded discounted tail permits a finite-horizon sign certificate. The distribution of individual gains remains unidentified; a finite transportation program gives sharp bounds and a two-point example computes them exactly. Network transfer accounting shows why tolls and capitalization cannot simply be added to gross benefits.
\end{abstract}
\section{A fixed project population and value convention}
Fix a cohort of candidate projects and a discount $\delta\in(0,1)$. Let $S_{jt}(a)$ be the gross monetary social service value at date $t$ under assignment $a\in\{0,1\}$, already net of congestion and environmental damages included in the chosen valuation. A fixed welfare population and network boundary are required. If other projects affect the site, either they are fixed as part of the intervention or assignment is defined at a sufficiently broad network level.

Let incremental maintenance under implementation be $M_{jt}$, construction outlay be $K_j$, and additional financing distortion or displaced-expenditure opportunity cost be $C_j^F$. Costs are in the same numeraire. Financing repayment of an already counted capital outlay is not counted again. For clarity first regard $K_j,C_j^F$, and the maintenance cost model as declared inputs. Define
\[
 V_j=\sum_{t\ge0}\delta^t
       \{S_{jt}(1)-S_{jt}(0)-M_{jt}\}-K_j-C_j^F.
 \tag{1}
\]
A completion delay is represented by the service stream, which is zero before service begins. Failed projects may produce zero services while retaining construction costs. Assume absolute integrability of all discounted sums.

The main statistical target is $\E V_j$ for the baseline candidate cohort. A target restricted to completed projects is different and can be selected by the treatment itself.
\section{Identification of the mean}
Let project assignment $Z_j$ be randomized with known propensity $e_j\in(0,1)$ independently of its potential service streams, conditional on any declared baseline strata. Put
$B_j(a)=\sum_t\delta^tS_{jt}(a)$ and
$C_j=K_j+C_j^F+\sum_t\delta^tM_{jt}$.
Assume $C_j$ is known for each candidate in this benchmark and $B_j^{\rm obs}=B_j(Z_j)$ is measured consistently. Then
\[
 \widehat V=\frac1n\sum_{j=1}^n
 \left[
 \frac{Z_j B_j^{\rm obs}}{e_j}
 -\frac{(1-Z_j)B_j^{\rm obs}}{1-e_j}
 -C_j
 \right]
 \tag{2}
\]
is unbiased for the finite-cohort mean of (1).

The proof conditions on the potential outcomes and known costs: the expected first weighted term is $B_j(1)$, and the expected second term is $B_j(0)$. Summing establishes the claim. No pairing of the two potential outcomes for the same project is observed or needed for the mean.

If costs are stochastic and realized only under treatment, they can instead be included in the treated outcome and inverse-probability weighted along with it, provided assignment is independent and the target is that incremental treated cost. The simplified subtraction in (2) must not be used as though a selectively observed cost were known for everyone. Likewise, unmeasured financing distortions remain unmeasured even in a randomized infrastructure experiment.

A valid local instrument would generally identify a different affected population and requires its own exclusion and compliance assumptions. Equation (2) is not a proof for every funding formula or timing instrument.
\section{A finite-horizon sign certificate}
Suppose for every project and date the incremental net service flow
$G_{jt}=S_{jt}(1)-S_{jt}(0)-M_{jt}$ obeys $|G_{jt}|\le H$ almost surely. Let $V_j^{(T)}$ truncate (1) after date $T$. The geometric series gives
\[
 |V_j-V_j^{(T)}|
 \le\sum_{t=T+1}^\infty\delta^t H
 =\frac{H\delta^{T+1}}{1-\delta}.
 \tag{3}
\]
The same bound holds for the difference between their means by taking expectations.

Thus if a valid lower confidence bound for $\E V_j^{(T)}$ exceeds the right side of (3), it is also a positive-value certificate for the infinite-horizon mean at that confidence level. If the upper bound is below its negative, negative mean value is certified. Otherwise the observed horizon and maintained tail bound do not determine the sign.

For instance, with $\delta=0.9$, $H=2$, and $T=20$, the omitted absolute tail is at most $20(0.9)^{21}$, approximately $2.1884$. A truncated mean estimate of $1$ alone cannot certify positive infinite-horizon value. The tail bound is an assumption about future net services, not something automatically established by observing the first twenty years.
\section{Sharp bounds for the fraction with positive gains}
Even perfect randomization generally identifies only marginal distributions of potential benefits. Consider a finite-support specialization with fixed known net cost $c$, treated discounted benefits $B(1)\in\{a_r\}_{r=1}^R$ with probabilities $p_r$, and control benefits $B(0)\in\{b_s\}_{s=1}^S$ with probabilities $q_s$.

A compatible joint distribution is a nonnegative array $\gamma_{rs}$ satisfying
\[
 \sum_s\gamma_{rs}=p_r,\qquad
 \sum_r\gamma_{rs}=q_s.
 \tag{4}
\]
The sharp lower and upper positive-value fractions are the minimum and maximum of
\[
 \sum_{r,s}\gamma_{rs}\ind\{a_r-b_s>c\}
 \tag{5}
\]
over (4).
\begin{theorem}
The linear programs (4)--(5) attain their optima and exactly characterize the identified interval for $\Prob\{V>0\}$ in this finite-support model.
\end{theorem}
\begin{proof}
The feasible arrays form a nonempty compact polytope: the product coupling $p_rq_s$ is feasible, and every coordinate lies in $[0,1]$. The objective is linear, so both extrema are attained. Every admissible population induces an array satisfying (4). Conversely each feasible array defines a population distribution of the two potential benefits, and random assignment independent of that pair reproduces the observed marginal distributions. Thus every feasible objective value is compatible. Convex mixtures of arrays fill the interval between the extrema.
\end{proof}
For a worked example, take $c=0$, $B(1)\in\{0,1\}$ with probability $0.6$ of one, and $B(0)\in\{0,1\}$ with probability $0.4$ of one. Let $x=\Prob(B(1)=1,B(0)=1)$. Nonnegative cell probabilities require $x\in[0,0.4]$. Positive gains occur exactly in cell $(1,0)$, whose probability is $0.6-x$. Hence the sharp positive-gain fraction is $[0.2,0.6]$, while mean gain is always $0.2$.

A positive average therefore does not identify the fraction of projects that benefit, much less which particular projects do so. Rank invariance would choose one coupling, but it is an additional restriction and is not supplied by treatment randomization.
\section{Network benefits and transfers}
Consider one trip with gross traveler benefit $b$, real travel resource cost $c$, toll $t$, and toll-operation resource cost $k$. Traveler surplus is $b-c-t$, and toll-operator surplus is $t-k$. Their sum is $b-c-k$. Treating the reduction of a toll as an additional resource saving after including both parties would double count a transfer. Unequal welfare weights can make the toll distributively relevant, but that is a different calculation.

Similarly, a permanent annual accessibility benefit can capitalize into land value. Adding both the entire discounted service benefit and the induced land-value increase to the same owners' welfare would duplicate the underlying stream. A project evaluation must choose an equivalent-variation or a gross-benefit convention and subtract each omitted real cost once. Route changes and relocation also mean that direct users are not necessarily the whole affected population.
\section{Remaining gap}
The attempt gives exact mean identification under randomization, a conditional tail certificate, and sharp finite-support distributional bounds. It does not measure project services, financing costs, spillovers, construction delays, or transportability in an actual infrastructure program. Out-of-sample prediction of positive net value needs baseline predictors, overlap, and validation beyond the mean identity.

These are standard evaluation and coupling arguments supplied with explicit assumptions. They do not resolve the broad public-investment research agenda or infer actual project returns from a spending-growth regression.

\section{Completion Estimate}
55\%

This is a post hoc, heuristic estimate for the full catalogue target.
The attempt identifies mean project value under randomization, certifies bounded discounted tails, and derives sharp finite-support positive-gain coupling bounds with consistent network accounting. Actual service and financing valuation, project spillovers and delays, heterogeneous returns, and out-of-sample transport still need evidence.

\begin{thebibliography}{9}
\bibitem{imf} International Monetary Fund (2025).
\emph{Fiscal Monitor, October 2025: Spending Smarter}.
The primary publisher summary concerns public-spending allocation and efficiency; it does not identify the project-specific counterfactuals imposed here.
\url{https://www.imf.org/en/publications/fm/issues/2025/10/07/fiscal-monitor-october-2025}.
\end{thebibliography}

\end{document}
